\chapter{Configuration \& Settings}
\label{ch:configuration}

\textit{Every tunable parameter in \sysname{} lives in a single file:
\codefile{apps/api/config.py}.  This chapter explains the settings architecture,
the security pattern for API keys, how to switch LLM and embedding providers,
and the rationale behind every cache TTL and rate limit.}


% ─────────────────────────────────────────────────────────────────────
\section{Settings Architecture}
\label{sec:settings-architecture}

\sysname{} uses Pydantic's \ic{BaseSettings} class from the
\ic{pydantic-settings} package.  Every configuration value is declared as a
typed field with a default, and values are automatically loaded from environment
variables or a \ic{.env} file.

\begin{lstlisting}[style=python, caption={The Settings class skeleton (simplified).}]
from pydantic import SecretStr
from pydantic_settings import BaseSettings, SettingsConfigDict

class Settings(BaseSettings):
    database_url: str = ""
    neo4j_uri: str = ""
    nvidia_api_key: SecretStr = SecretStr("")
    llm_provider: str = "nvidia"
    fuzzy_match_threshold: float = 0.85
    rate_limit_requests: int = 30
    debug: bool = False
    # ... 30+ more fields

    model_config = SettingsConfigDict(
        env_file=".env", extra="ignore"
    )
\end{lstlisting}

The \ic{model\_config} line tells Pydantic two things:

\begin{enumerate}
    \item Load variables from a \ic{.env} file in the working directory.
    \item Ignore extra environment variables that do not correspond to any
          field---this prevents crashes when the shell has unrelated variables set.
\end{enumerate}

The entire application accesses settings through a singleton factory:

\begin{lstlisting}[style=python, caption={The singleton pattern for settings access.}]
from functools import lru_cache

@lru_cache
def get_settings() -> Settings:
    return Settings()
\end{lstlisting}

The \ic{@lru\_cache} decorator ensures that \ic{Settings()} is constructed
exactly once.  Every subsequent call to \ic{get\_settings()} returns the same
instance with zero overhead.  This matters because Pydantic parses the
environment on every \ic{Settings()} call---without caching, that parsing would
happen on every request.


% ─────────────────────────────────────────────────────────────────────
\section{The SecretStr Pattern}
\label{sec:secretstr-pattern}

Three fields in \ic{Settings} use Pydantic's \ic{SecretStr} type instead of
plain \ic{str}:

\begin{itemize}
    \item \ic{nvidia\_api\_key} --- NVIDIA NIM API key
    \item \ic{neo4j\_password} --- Neo4j database password
    \item \ic{secret\_key} --- JWT signing secret
\end{itemize}

\ic{SecretStr} wraps the raw value so that it is never accidentally exposed.
Calling \ic{print(settings)} or logging the object shows
\ic{SecretStr('**********')} instead of the actual key.

\subsection*{Accessing the raw value}

You cannot use a \ic{SecretStr} field directly as a string.  The \ic{Settings}
class provides explicit accessor methods:

\begin{lstlisting}[style=python, caption={Correct vs.\ incorrect access of secret fields.}]
settings = get_settings()

# CORRECT: use the accessor method
key = settings.get_nvidia_api_key()     # returns str
pwd = settings.get_neo4j_password()     # returns str
sec = settings.get_secret_key()         # returns str

# ALSO CORRECT: call .get_secret_value() directly
key = settings.nvidia_api_key.get_secret_value()

# WRONG: this gives you a SecretStr object, not a string
key = settings.nvidia_api_key  # type: SecretStr, not str
\end{lstlisting}

\begin{keyinsight}
The \ic{SecretStr} pattern (tagged SEC-007 in the codebase) is a defense-in-depth
measure.  Even if a developer accidentally writes
\ic{logger.info(f"Config: \{settings\}")}, the API keys remain masked.
The custom \ic{\_\_repr\_\_} method on \ic{Settings} adds a second layer by only
including non-sensitive fields in the string representation.
\end{keyinsight}


% ─────────────────────────────────────────────────────────────────────
\section{LLM Provider Switching}
\label{sec:llm-provider-switching}

\sysname{} supports two LLM providers, selected by a single environment
variable:

\begin{lstlisting}[style=bash, numbers=none]
# In your .env file:
LLM_PROVIDER=nvidia    # default
# or
LLM_PROVIDER=bedrock   # Amazon Bedrock
\end{lstlisting}

The application uses a factory pattern to instantiate the correct provider at
startup.  The provider interface is identical regardless of backend, so no
application code needs to change when you switch.

\subsection*{NVIDIA NIM (default)}

\begin{itemize}
    \item \textbf{Primary model:}
          \ic{nvidia/llama-3.3-nemotron-super-49b-v1.5} (128k context)
    \item \textbf{Fallback model:}
          \ic{nvidia/llama-3.1-nemotron-70b-instruct} (auto-failover on 503 errors)
    \item Requires \confkey{NVIDIA\_API\_KEY} from \ic{build.nvidia.com}
\end{itemize}

The fallback is automatic.  When the primary model returns a 503 (common under
load on the free NIM tier), the system retries with the fallback model.  This
is configured via:

\begin{lstlisting}[style=python, numbers=none]
llm_fallback_model: str = "nvidia/llama-3.1-nemotron-70b-instruct"
llm_fallback_enabled: bool = True
\end{lstlisting}

\subsection*{Amazon Bedrock (alternative)}

\begin{itemize}
    \item \textbf{Model:} \ic{anthropic.claude-sonnet-4-20250514}
    \item Requires AWS credentials in the environment (no explicit API key field).
    \item Region defaults to \ic{us-west-2}, configurable via
          \confkey{AWS\_REGION}.
\end{itemize}


% ─────────────────────────────────────────────────────────────────────
\section{Embedding Provider}
\label{sec:embedding-provider}

Embeddings use a separate provider selection:

\begin{lstlisting}[style=bash, numbers=none]
EMBEDDING_PROVIDER=nvidia    # default
\end{lstlisting}

The default model is \ic{nvidia/llama-3.2-nv-embedqa-1b-v2}, which produces
2048-dimensional vectors.  A separate API key (\confkey{NVIDIA\_EMBEDDING\_API\_KEY})
is required because NVIDIA issues distinct keys for LLM and embedding endpoints.

\begin{keyinsight}
You can use NVIDIA embeddings with the Bedrock LLM, and this is the recommended
configuration if you switch LLM providers.  Set \ic{LLM\_PROVIDER=bedrock} but
keep \ic{EMBEDDING\_PROVIDER=nvidia} to avoid re-indexing.
\end{keyinsight}

\begin{warning}
Changing the embedding provider changes the vector dimensionality (NVIDIA
produces 2048-d, Bedrock Titan produces 1024-d).  Existing vectors in Neo4j
become incompatible---you must delete all \ic{embedding} properties and re-embed
every decision trace and entity.  This is a multi-hour operation on a large
graph.  Do not change the embedding provider casually.
\end{warning}


% ─────────────────────────────────────────────────────────────────────
\section{Cache TTL Rationale}
\label{sec:cache-ttl-rationale}

Every cache TTL in \sysname{} was chosen based on how quickly the underlying
data can change.  Table~\ref{tab:cache-ttl-rationale} explains the reasoning.

\begin{table}[ht]
\centering
\caption{Cache TTL rationale.}
\label{tab:cache-ttl-rationale}
\begin{tabularx}{\textwidth}{l c X}
\toprule
\textbf{Cache} & \textbf{TTL} & \textbf{Why This Value} \\
\midrule
Entity resolution & 5 min &
    Entities can change when merged or when aliases are added.
    A short TTL ensures stale resolutions expire quickly.
    Even 5 minutes captures the burst of repeated mentions
    within a single extraction. \\
\addlinespace
LLM response & 24 h &
    The same input text with the same prompt version produces the same
    extraction.  A 24-hour TTL is safe because the cache is
    keyed by prompt version---bumping
    \confkey{LLM\_EXTRACTION\_PROMPT\_VERSION} invalidates all entries. \\
\addlinespace
Embedding vector & 30 d &
    The embedding model does not change between deployments.
    The same text always maps to the same vector.  A 30-day TTL
    minimizes redundant API calls while still eventually refreshing
    if the model is swapped. \\
\bottomrule
\end{tabularx}
\end{table}


% ─────────────────────────────────────────────────────────────────────
\section{Rate Limiting}
\label{sec:rate-limiting}

\sysname{} implements per-user rate limiting via a Redis-backed token bucket.
The parameters:

\begin{itemize}
    \item \textbf{Authenticated users:} 30 requests per 60-second window
    \item \textbf{Anonymous users:} 10 requests per 60-second window
    \item \textbf{Window type:} Sliding window (counter resets 60 seconds after
          the first request in the current window)
\end{itemize}

The Redis key follows the pattern \ic{rate:\{user\_id\}} with a 60-second TTL.
When the counter exceeds the limit, the API returns HTTP 429 (Too Many Requests)
with a \ic{Retry-After} header.

Both parameters are configurable:

\begin{lstlisting}[style=python, numbers=none]
rate_limit_requests: int = 30   # requests per window
rate_limit_window: int = 60     # window size in seconds
\end{lstlisting}


% ─────────────────────────────────────────────────────────────────────
\section{Key Settings Reference}
\label{sec:key-settings-reference}

Table~\ref{tab:key-settings} lists the 15 most important settings.
Appendix~\ref{ch:config-reference} contains the complete list.

\begin{longtable}{l l p{6.5cm}}
\caption{Key settings in \codefile{config.py}.}
\label{tab:key-settings} \\
\toprule
\textbf{Setting} & \textbf{Default} & \textbf{Description} \\
\midrule
\endfirsthead
\toprule
\textbf{Setting} & \textbf{Default} & \textbf{Description} \\
\midrule
\endhead
\midrule
\multicolumn{3}{r}{\itshape Continued on next page\ldots} \\
\bottomrule
\endfoot
\bottomrule
\endlastfoot
\confkey{LLM\_PROVIDER}
    & \ic{nvidia}
    & LLM backend: \ic{nvidia} or \ic{bedrock}. \\
\confkey{EMBEDDING\_PROVIDER}
    & \ic{nvidia}
    & Embedding backend. Keep \ic{nvidia} to avoid re-indexing. \\
\confkey{NVIDIA\_MODEL}
    & \ic{nvidia/llama-3.3-nemotron-super-49b-v1.5}
    & Primary LLM model identifier. \\
\confkey{LLM\_FALLBACK\_MODEL}
    & \ic{nvidia/llama-3.1-nemotron-70b-instruct}
    & Fallback LLM on 503 errors. \\
\confkey{LLM\_FALLBACK\_ENABLED}
    & \ic{True}
    & Enable automatic model fallback. \\
\confkey{FUZZY\_MATCH\_THRESHOLD}
    & \ic{0.85}
    & Minimum fuzzy match score (0--1) for entity resolution. \\
\confkey{EMBEDDING\_SIMILARITY\_THRESHOLD}
    & \ic{0.90}
    & Minimum cosine similarity for embedding-based entity match. \\
\confkey{ENTITY\_CACHE\_TTL}
    & \ic{300}
    & Entity cache lifetime in seconds (5 minutes). \\
\confkey{LLM\_CACHE\_TTL}
    & \ic{86400}
    & LLM response cache lifetime in seconds (24 hours). \\
\confkey{LLM\_EXTRACTION\_PROMPT\_VERSION}
    & \ic{v1}
    & Bump to invalidate all LLM cache entries. \\
\confkey{MAX\_PROMPT\_TOKENS}
    & \ic{70000}
    & Maximum input tokens for LLM calls (128k context). \\
\confkey{RATE\_LIMIT\_REQUESTS}
    & \ic{30}
    & Requests per rate limit window. \\
\confkey{RATE\_LIMIT\_WINDOW}
    & \ic{60}
    & Rate limit window in seconds. \\
\confkey{SECRET\_KEY}
    & (empty)
    & JWT signing secret. Must match \confkey{NEXTAUTH\_SECRET}. \\
\confkey{DEBUG}
    & \ic{False}
    & Enable debug mode (verbose logging, CORS relaxed). \\
\end{longtable}
